\documentclass[12pt]{article}
\setlength{\parskip}{5mm}
\usepackage[T1]{fontenc}
\usepackage[utf8]{inputenc}
\usepackage{url, hyperref}
\usepackage[top = 2.5 cm, bottom = 2.5 cm, left = 1.7cm, right = 1.7cm]{geometry}
\usepackage{amsmath, amsthm, amsfonts, amssymb, enumitem, float}
\usepackage[spanish]{babel}
\usepackage{graphicx}
\graphicspath{ {./} }
\usepackage{fancyhdr}
\usepackage{tikz}
\usepackage{listings}
\usepackage{xcolor}
\usetikzlibrary{arrows,calc}
\tikzstyle{edge}=[shorten <=2pt, shorten >=2pt, >=stealth, line width=1.1pt]
\tikzstyle{vertex}=[circle, draw, minimum size=5pt, inner sep=0pt,
outer sep=0pt,line width=0.7pt]

\pagestyle{fancy}
\setlength{\headheight}{15pt}
\fancyhf{}

\newcommand{\QED}{\begin{flushright}$\blacksquare$\end{flushright}}

\lstset{
  language=Python,        % lenguaje del código
  basicstyle=\ttfamily,   % estilo de la fuente
  keywordstyle=\color{blue},
  stringstyle=\color{red},
  commentstyle=\color{green!60!black},
%   numbers=left,           % numerar líneas
  numberstyle=\tiny,
  stepnumber=1,
  showstringspaces=false,
  inputencoding=utf8,  % para que lea UTF-8
  extendedchars=true,  % permite algunos caracteres extendidos
  literate={á}{{\'a}}1 {é}{{\'e}}1 {í}{{\'i}}1 {ó}{{\'o}}1 {ú}{{\'u}}1 
           {Á}{{\'A}}1 {É}{{\'E}}1 {Í}{{\'I}}1 {Ó}{{\'O}}1 {Ú}{{\'U}}1 
           {ñ}{{\~n}}1 {Ñ}{{\~N}}1 {π}{{$\pi$}}1 {≈}{{$\approx$}}1}

\title{Extracción y perfilado de datos}

\begin{document}

\lhead{Extracción y perfilado de datos}

\title{Práctica 4 - Extracción y perfilado de datos}

\thispagestyle{empty}

\begin{figure}[ht]
	\minipage{0.75\textwidth}
	\includegraphics[width = 4.6 cm]{logo_unam.png}
	\endminipage
	\minipage{0.32\textwidth}
	\includegraphics[height = 4.9cm, width = 4cm]{logo_iimas.png}
	\endminipage 
\end{figure}

\begin{center}
	\vspace{0.5cm}
	\LARGE
	UNIVERSIDAD NACIONAL AUTÓNOMA DE MÉXICO

	\vspace{0.5cm}
	\LARGE
	Instituto de Investigaciones en Matemáticas Aplicadas y Sistemas

	\vspace{1.5cm}
	\Large
	\textbf{Práctica 4 DWH - Extracción y perfilado de datos}

	\vspace{1cm}
	\normalsize
	PRESENTA \\
	\vspace{.3cm}
	\large
	\textbf{Díaz Juárez Ana Sofía \\}
	\textbf{Munive Ramírez Ibrahim \\}
	\textbf{Pérez Aguiar Oropeza Gabriel Emiliano \\}


	\vspace{1cm}
	\normalsize
	PROFESORA \\
	\vspace{.3cm}
	\large
	\textbf{Dra. María del Pilar Angeles}

	\vspace{1cm}
	\normalsize
	ASIGNATURA \\
	\vspace{.3cm}
	\large
	\textbf{Bases de Datos Estructuradas}

	\vspace{1cm}
\end{center}

\newpage


\section*{Contexto}

Una vez creado el repositorio ServiciosFinancieros, es necesario
explotar las funetes OLTP con las que se cuenta. Considerar el 
archivo \texttt{database-sucia-citi.csv} que contiene las 
direcciones de las sucursales de Citi Bank que servirá para cargar 
datos limpios a la tabla Branch.

\begin{center}
    \begin{tabular}{|c|}
        \hline
        \underline{\textbf{branch\_key}} \\
        branch\_name \\
        branch\_addres \\
        branch\_city \\
        branch\_state \\
        branch\_zip \\
        branch\_type \\
        \hline
    \end{tabular}
\end{center}

\vspace{0.5 cm}
\hrule
\vspace{0.5 cm}

\section*{Objetivos}

Extraer los datos relevantes de la fuente dada y realizar perfilado 
de calidad de datos para el sistema de Servicios Financieros. 


\vspace{1.5 cm}
\hrule
\vspace{0.5 cm}

A partir del archivo \texttt{database-sucia-citi.csv} extraer 
aquellos datos que son relevanres para ser cargados a la tabla 
Branch del repositorio ServiciosFinancieros y guardar el archivo 
'Branch.csv' con campos delimitados por comas con nombre.



Obtener el perfil de los datos almacenados en \texttt{Branch.csv} 
siguiendo el procedimiento

\begin{enumerate}
    \item Explorar los datos contenidos en \texttt{Branch.csv}
    
    \item Idenfificar y anotar las características que tiene 
        cada campo, por ejemplo: si tiene valores faltances, 
        falta de homogeneidad en el formato de los textos, etc.
    
    \item Generar las siguientes estadísticas utilizando el 
        software de su preferencia:
        \begin{itemize}
            \item[$\mathcal{A}1$.] Número de registros del archivo
            \item[$\mathcal{A}2$.] Número de valores faltantes
            \item[$\mathcal{A}3$.] Número de registros duplicados
            \item[$\mathcal{A}4$.] Número de códigos postales de 
                Disparos Unidos de América válidos     
        \end{itemize}
    
    \item Identificar otros indicadores estadísticos que puedan 
        servir para el análisis OLAP
    
    \item Identificar si existe información que se requiera en la 
        tabla Branch que no la proporcione \texttt{database-sucia-citi.csv}
        y explciar cómo completar dicha información. 

\end{enumerate}

\QED

\end{document}
